\chapter{Transformada de Fourier de tiempo continuo}

\section{Introducci\'{o}n}

\begin{enumerate}

\item   Vamos a ver como podemos \textbf{modificar} el concepto de \textbf{serie de Fourier} \textbf{para estudiar funciones no peri\'{o}dicas} de tiempo continuo.  
Esto nos permite \textbf{definir} el concepto de \textbf{contenido de frecuencias de una se\~{n}al}, a\'{u}n cuando sea no peri\'{o}dica.

\item   Pensemos en el ejemplo (\ref{funcion.par}), donde
\begin{equation}
	a_{0} = \frac{T_{1}}{T_{2}},
\end{equation}
y
\begin{equation}
	a_{k} =  \frac{T_{1}}{T} \, \frac{sin(\frac{k \, \pi \, T_{1}}{T})}{\frac{k \, \pi \, T_{1}}{T}}
	. \label{tf.1}
\end{equation}
\item  Podemos escribir (\ref{tf.1}) de la siguiente forma
        \begin{equation}
           T a_{k}=
	T \, a_{k} =  \left. T_{1} \, \frac{sin(\omega)}{\omega} \right|_{\omega = \frac{k \, \pi \, T_{1}}{T}}
        \end{equation}
        y entonces es evidente que la podemos interpretar como una \textbf{secuencia} formada por \textbf{muestras} de una \textbf{funci\'{o}n envolvente} $T_{1} \, \frac{\sin(\omega)}{\omega}$) donde la variable $\omega $ es continua, y 
        la secuencia $T a_{k}$ está formada por valores igualmente espaciados en 
        frecuencia, con intervalo $\frac{\pi \, T_{1}}{T}$.
        
        \item   Observe que la separaci\'{o}n de las muestras de la envolvente de $T a_{k}$ est\'{a} dada 
        por 
        \begin{equation}
                \omega_{T} = \frac{\pi \, T_{1}}{T}.
        \end{equation}
        Entonces, cuando mayor sea $T$ menor ser\'{a} $\omega_{T}$ y la cantidad de muestras en un intervalo dado de frecuencia ser\'{a} mayor.

        \item  Si fijamos el valor $T_{1}$ (recordando que debe ser $T_{1}<T,$ por hip\'{o}tesis) la envolvente de $T a_{k}$ es independiente de $T$.
\end{enumerate}

\section{Síntesis y análisis}

\subsection{Teor\'{\i}a}

\begin{enumerate}

        \item  Vamos a construir una se\~{n}al $x\left( t\right) $ no peri\'{o}dica como el l\'{\i}mite de 
        una se\~{n}al $\hat{x}\left( t\right) $ peri\'{o}dica cuando el per\'{\i}odo $T$ tiende a 
        $\infty .$ Veamos que pasa con los coeficientes de la serie de Fourier. 
        Recordemos que una se\~{n}al peri\'{o}dica $\hat{x}\left( t\right) =\hat{x}\left( t+T\right) $ 
        se puede escribir como una serie de Fourier mediante la ecuaci\'{o}n de s\'{\i}ntesis 
        \begin{equation}
                \hat{x}\left( t\right) =\sum_{k=-\infty }^{\infty }a_{k} \, e^{j \, k \, \omega_{0} \, t},
        \end{equation}
        donde los coeficientes $a_{k}$ se calculan mediante la ecuaci\'{o}n de an\'{a}lisis 
        \begin{equation}
            a_{k}=\frac{1}{T}\int_{-\frac{T}{2}}^{\frac{T}{2}}\hat{x}\left( t\right) e^{-j \, k \, \omega_{0} \, t}dt.
        \end{equation}

        \item  Definamos la se\~{n}al no peri\'{o}dica $x(t)$ especificando 
        \begin{equation}
                x\left( t\right) =
                        \left\{ \begin{array}{ll}
                                \hat{x}\left( t\right) & \left| t\right| \leq \frac{T}{2}, \\
                                & \\
                                0 & \left| t\right| >\frac{T}{2},
                                \end{array} 
                        \right. \qquad
                        \hat{x}\left( t\right) =\hat{x}\left( t+T\right) .
        \end{equation}
        Entonces los coeficientes de la serie de Fourier son 
        \begin{equation}
           a_{k}=
                  \frac{1}{T}\int_{-\frac{T}{2}}^{\frac{T}{2}}\hat{x}\left( t\right)e^{-j \, k \, \omega_{0} \, t}dt=
                  \frac{1}{T}\int_{-\infty }^{\infty }x\left( t\right)e^{-j \, k \, \omega_{0} \, t}dt.
        \end{equation}
        \textbf{Las dos integrales son iguales en valor porque un per\'{\i}odo de }
        $\hat{x}\left( t\right) $  \textbf{\ coincide con }$x\left( t\right) ,$ 
        \textbf{\ justamente donde se eval\'{u}a la primer integral, y } $x\left(t\right) $
        \textbf{\ es nula para }$\left| t\right| >\frac{T}{2}.$

        \item  Definamos la ecuaci\'{o}n de an\'{a}lisis de la transformada de Fourier como 
        \begin{equation}
                X\left( j \, \omega \right) =\int_{-\infty }^{\infty }x\left( t\right)e^{-j \, \omega \, t} \, dt.
        \end{equation}
        Tenemos que los coeficientes de la serie de Fourier forman la secuencia de valores 
        \begin{equation}
                Ta_{k}=X\left( j \, k \, \omega_{0}\right) .
        \end{equation}

        \item  Usando la ecuaci\'{o}n de s\'{\i}ntesis de la serie de Fourier podemos escribir 
        \begin{equation}
            \hat{x}\left( t\right) =
               \frac{1}{T}\sum_{k=-\infty }^{\infty }X\left(j \, k \, \omega_{0}\right) e^{j \, k \, \omega_{0} \, t}.
        \end{equation}
        Dado que $\omega _{0}=\frac{2 \, \pi }{T}\longleftrightarrow T=\frac{2 \, \pi }{\omega _{0}}$ 
        vemos que la serie de Fourier puede escribirse como
        \begin{equation}
           \hat{x}\left( t\right) =
           \frac{1}{2 \, \pi }
              \sum_{k=-\infty }^{\infty }X\left(j \, k \, \omega_{0}\right) e^{j \, k \, \omega_{0} \, t} \, \omega _{0}.
        \end{equation}
        A medida que $T\rightarrow \infty ,$ $\omega _{0}\rightarrow 0$ (o sea un diferencial de 
        frecuencia $d\omega $), $\hat{x}\left( t\right) \rightarrow x\left( t\right) $ y obtenemos la 
        ecuaci\'{o}n de s\'{\i}ntesis 
        \begin{equation}
           x\left( t\right) =
              \frac{1}{2 \, \pi }\int_{-\infty }^{\infty }X\left( j \, \omega \right) e^{j \, \omega \, t}d\omega .
        \end{equation}
        \textbf{A medida que }$T\rightarrow \infty ,$\textbf{\ se agranda el rango de valores de } $t$ 
        \textbf{\ en la que las dos funciones }$\hat{x}\left(t\right) $\textbf{\ e }$x\left( t\right) $
        \textbf{\ son iguales.}
\end{enumerate}

\subsection{Prestar atenci\'{o}n}

Como hemos visto antes, vamos a usar la ecuaci\'{o}n de an\'{a}lisis 
\begin{equation}
        X\left( j \, \omega \right) =
                \int_{-\infty }^{\infty }x\left( t\right) e^{-j \, \omega \, t} \, dt,  \label{analisis}
\end{equation}
y la ecuaci\'{o}n de s\'{\i}ntesis 
\begin{equation}
    x\left( t\right) =
   \frac{1}{2 \, \pi }\int_{-\infty }^{\infty }X\left( j \, \omega \right) e^{j \, \omega \, t}d\omega .  \label{sintesis}
\end{equation}
Esto significa que estamos usando la variable $\omega $ (radianes por segundo), y no la variable $f$ 
(ciclos por segundo). Recordemos que ambas est\'{a}n relacionadas por 
\begin{equation}
        \omega =2 \, \pi f.
\end{equation}
Este cambio de variables afecta a las ecuaciones que ser\'{\i}an de la forma 
\begin{equation}
        X\left( f\right) =\int_{-\infty }^{\infty }x\left( t\right) e^{-j2 \, \pi ft}dt,
\end{equation}
y 
\begin{equation}
        x\left( t\right) =\int_{-\infty }^{\infty }X\left( f\right) e^{j2 \, \pi ft}df.
\end{equation}
Usamos las f\'{o}rmulas dadas por las ecs. (\ref{analisis}) y (\ref{sintesis}) porque son las que aparecen 
en los libros disponibles en biblioteca. La diferencia de notaci\'{o}n no afecta la teor\'{\i}a.

\subsection{Concepto de la Transformada de Fourier}

\begin{enumerate}

        \item   La Transformada de Fourier es una generalizaci\'{o}n de la serie de Fourier, y permite estudiar funciones no peri\'{o}dicas en el dominio de la 
        frecuencia.
        
        \item   La Transformada de Fourier permite descomponer una se\~{n}al en componentes arm\'{o}nicas cuyas frecuencias varían en diferenciales, a diferencia de la serie de Fourier donde las frecuencias de las arm\'{o}nicas varían en incrementos de $\omega_{T} = \frac{2 \, \pi}{T}$ la frecuencia fundamental correspondiente al período $T$.

\end{enumerate}

\section{Convergencia de la Transformada de Fourier}

Sea una se\~{n}al $x\left( t\right) ,$ y consideremos las funciones $X\left(j \, \omega \right) $ y 
$x_{F}\left( t\right) $ definidas como 
\begin{equation}
X\left( j \, \omega \right) =
        \int_{-\infty }^{\infty }x\left( t\right) e^{-j \, \omega \, t} \, dt,  \label{convergenciaUno}
\end{equation}
y 
\begin{equation}
x_{F}\left( t\right) =
        \frac{1}{2 \, \pi }\int_{-\infty }^{\infty }X\left(j \, \omega \right) e^{j \, \omega \, t}d\omega .
\end{equation}
La pregunta es: \textquestiondown cu\'{a}ndo $x_{F}\left( t\right) $ es una representaci\'{o}n 
v\'{a}lida de $x\left( t\right) ?$

\subsection{Energ\'{\i}a finita}

Si $x\left( t\right) $ tiene energ\'{\i}a finita 
\begin{equation}
        \int_{-\infty }^{\infty }\left[ x\left( t\right) \right] ^{2}dt<\infty ,
\end{equation}
entonces $X\left( j \, \omega \right) $ es finita, o sea que la integral (\ref{convergenciaUno}) converge, y 
adem\'{a}s 
\begin{equation}
        \int_{-\infty }^{\infty }\left[ x\left( t\right) -x_{F}\left( t\right) \right] ^{2}dt=0.
\end{equation}
Esto implica que $x_{F}\left( t\right) $ y $x\left( t\right)$ son similares excepto en, a lo sumo, una cantidad limitada de valores de $t$ donde las funciones tengan discontinuidades.

\subsection{Condiciones de Dirichlet}

Son \textbf{condiciones alternativas a la anterior}, y son \textbf{suficientes} para asegurar que 
$x_{F}\left( t\right) $ es igual a $x\left(t\right) $ excepto en puntos de discontinuidad.

\begin{enumerate}
        \item  $x\left( t\right) $ debe ser absolutamente integrable 
        \begin{equation}
                \int_{-\infty }^{\infty }\left| x\left( t\right) \right| dt<\infty .
        \end{equation}

        \item  $x\left( t\right) $ tiene un n\'{u}mero finito de m\'{a}ximos y m\'{\i}nimos dentro de 
        cualquier intervalo finito.

        \item  $x\left( t\right) $ tiene un n\'{u}mero finito de discontinuidades dentro de cualquier 
        intervalo finito, y dichas discontinuidades son finitas. En este caso, en los puntos de 
        discontinuidad de $x\left( t\right)$, $x_{F}\left( t\right)$ es 
        igual al promedio de valores de $x\left( t\right)$ a ambos lados de la discontinuidad.
\end{enumerate}
Si se cumplen estas condiciones entonces $x\left( t\right) $ tiene transformada de Fourier.

\subsection{Pregunta}

Las se\~{n}ales que llegan a las antenas de un televisor o un equipo de radio, por mencionar dos ejemplos, 
\textquestiondown tienen transformada de Fourier? Justifique.

\section{Ejemplos}



\subsection{Impulso temporal}

Calculemos la transformada de Fourier de un impulso $\delta \left( t\right) $
\begin{equation}
        X\left( j \, \omega \right) =  \int_{-\infty }^{\infty }\delta \left( t\right)e^{-j \, \omega \, t} \, dt.
\end{equation}
Por la propiedad del muestreo 
\begin{equation}
\delta \left( t\right) e^{-j \, \omega \, t}=
        \delta \left( t\right) e^{-j \, \omega0}=\delta \left( t\right) ,
\end{equation}
y por lo tanto 
\begin{equation}
X\left( j \, \omega \right) =
        \int_{-\infty }^{\infty }\delta \left( t\right)e^{-j \, \omega \, t} \, dt=
           \int_{-\infty }^{\infty }\delta \left( t\right) dt=1.
\end{equation}
Note que es un valor constante para todo valor de $\omega .$

\subsection{Impulso temporal desplazado}

Calculemos la transformada de Fourier de un impulso desplazado en el 
tiempo $\delta \left( t - t_{0}\right) $
\begin{equation}
        X\left( j \, \omega \right) =
                \int_{-\infty }^{\infty }\delta \left( t - t_{0} \right)e^{-j \, \omega \, t} \, dt.
\end{equation}
Recordemos que, por la propiedad del muestreo,
\begin{equation}
\delta \left( t - t_{0} \right) e^{-j \, \omega \, t}=
        \delta \left( t - t_{0} \right) e^{-j \, \omega \, t_{0}} =
                \delta \left( t - t_{0} \right) e^{-j \omega \, t_{0}}
\end{equation}
y por lo tanto  podemos escribir
\begin{equation}
X\left( j \, \omega \right) =
        \int_{-\infty }^{\infty }\delta \left( t - t_{0} \right)e^{-j \, \omega \, t} \, dt=
        \left(\int_{-\infty }^{\infty }\delta \left( t\right) dt\right) \, e^{-j \, \omega \, t_{0}} = e^{-j \, \omega \, t_{0}},
\end{equation}
teniendo en cuenta que el factor $e^{-j \, \omega \, t_{0}}$ no depende de la variable de integración $t$ y puede ser extraído de la integral como una constante.


\subsection{Impulso en frecuencia}

Sea la transformada de Fourier 
\begin{equation}
        X\left( j \, \omega \right) =2 \, \pi \, \delta \left( \omega -\omega _{0}\right) .
\end{equation}
Apliquemos la ecuaci\'{o}n de s\'{\i}ntesis para obtener 
\begin{equation}
        x\left( t\right) =
          \frac{1}{2 \, \pi }
           \int_{-\infty }^{\infty }
               2 \, \pi \, \delta \left(\omega -\omega _{0}\right) \, e^{j \, \omega \, t} \, d\omega = e^{j \, \omega _{0} \, t}.
\end{equation}
Esto significa que las funciones periódicas $e^{j \, \omega _{0} \, t}$ también tienen transformada de Fourier, y son impulsos en frecuencia $2 \, \pi \, \delta \left( \omega -\omega _{0}\right)$.

\subsection{Tren de impulsos en frecuencia}

En general, si tenemos un tren de impulsos en frecuencia
\begin{equation}
        X\left( j \, \omega \right) =
                \sum_{k=-\infty }^{\infty }2 \, \pi \, a_{k} \, \delta \left(\omega -k \, \omega_{0}\right) ,
\end{equation}
y la serie de Fourier
\begin{equation}
        x\left( t\right) =\sum_{k=-\infty }^{\infty }a_{k} \, e^{j \, k \, \omega_{0} \, t},
\end{equation}
podemos escribir 
\begin{equation}
        \sum_{k=-\infty }^{\infty }a_{k} \, e^{j \, k \, \omega_{0} \, t} 
        \begin{array}{c}
                TF \\ 
                \Leftrightarrow
        \end{array}
        \sum_{k=-\infty }^{\infty }2 \, \pi \, a_{k} \, \delta \left( \omega -k \omega_{0} \right). \label{tf.tif}
\end{equation}

\subsection{Tren de impulsos temporal}

Consideremos el tren de impulsos temporal
\begin{equation}
        x\left( t\right) =\sum_{k=-\infty }^{\infty }\delta \left( t-kT\right) .
\end{equation}
Calculemos los coeficientes de la serie de Fourier, y observemos que dentro de los l\'{\i}mites de 
integraci\'{o}n $x\left( t\right) = \delta \left(t\right) $ lo que permite escribir 
\begin{equation}
        a_{k}=
        \frac{1}{T}\int_{-\frac{T}{2}}^{\frac{T}{2}}\delta \left( t\right)e^{-j \, k \, \omega \, t} \, dt=\frac{1}{T}.
\end{equation}
y por lo tanto tenemos que la transformada de Fourier es (vea \ref{tf.tif})
\begin{equation}
        X\left( j \, \omega \right) =
                \sum_{k=-\infty }^{\infty }\frac{2 \, \pi}{T} \, \delta\left( \omega -\frac{2 \, \pi \, k}{T}\right).
\end{equation}

\subsection{Pulso rectangular en tiempo}

Calculemos la transformada de la se\~{n}al pulso rectangular 
\begin{equation}
x\left( t\right) =
        \left\{ \begin{array}{l}
                1, \left| t\right| \leq T_{1}, \\ 
                0, \left| t\right| > T_{1}.
        \end{array} \right.
\end{equation}
Seg\'{u}n la ecuaci\'{o}n de an\'{a}lisis 
\begin{equation}
        X\left( j \, \omega \right) =
                \int_{-T_{1}}^{T_{1}}e^{-j \, \omega \, t} \, dt=
                \frac{\left[e^{-j \, \omega \, t}\right]_{-T_{1}^{T_{1}}}}{-j \, \omega}.
\end{equation}
Evaluando la integral
\begin{equation}
	X\left( j \, \omega \right) =
	\frac{\left[e^{-j \, \omega \, t}\right]_{-T_{1}}^{T_{1}}}{-j \, \omega} =
	\frac{e^{-j \, \omega \, T_{1}}-e^{j \, \omega \, T_{1}}}{-j \, \omega}.
\end{equation}
Finalmente, multiplicando y diviendo por $2 \, T_{1}$ y simplificando signos llegamos a la expresión
\begin{equation}
	X\left( j \, \omega \right) =
	\frac{e^{-j \, \omega \, T_{1}}-e^{j \, \omega \, T_{1}}}{-j \, \omega} =	
    2 \, T_{1}	\, \frac{e^{j \, \omega \, T_{1}}-e^{-j \, \omega \, T_{1}}}{2 \, j \, \omega \, T_{1}} =	
	2 \, T_{1} \, \frac{\sin\left( \omega \, T_{1}\right) }{\omega \, T_{1}}.
\end{equation}
Por lo tanto, resulta
\begin{equation}
	X\left( j \, \omega \right) =
	2 \, T_{1} \, \frac{\sin\left( \omega \, T_{1}\right) }{\omega \, T_{1}}.
\end{equation}

\subsection{Pulso rectangular en frecuencia}

Veamos cu\'{a}l es la se\~{n}al $x\left( t\right) $ cuya transformada de Fourier es 
\begin{equation}
        X\left( j \, \omega \right) =
                \left\{ \begin{array}{l}
                        1, \left| \omega \right| \leq W_{1}, \\ 
                        0, \left| \omega \right| >W_{1}.
                \end{array} \right.
\end{equation}
Seg\'{u}n la ecuaci\'{o}n de s\'{\i}ntesis 
\begin{equation}
        x\left( t\right) =
                \frac{1}{2 \, \pi}\int_{-W_{1}}^{W_{1}} e^{j \, \omega \, t}d\omega =
                        \frac{\sin \left( W_{1}t\right) }{\pi t},
\end{equation}
Evaluando la integral
\begin{equation}
	x\left( t \right) =
	\frac{1}{2 \, \pi} \, \frac{\left[e^{-j \, \omega \, t}\right]_{-W_{1}}^{W_{1}}}{-j \, t} =
	\frac{1}{2 \, \pi} \, \frac{e^{-j \,  W_{1} \, t}-e^{j \, W_{1} \, t}}{-j \, t}.
\end{equation}
Finalmente, multiplicando y diviendo por $2 \, W_{1}$ y simplificando signos llegamos a la expresión
\begin{equation}
	x\left( t \right) =
	\frac{1}{2 \, \pi} \, \frac{e^{-j \,  W_{1} \, t}-e^{j \, W_{1} \, t}}{-j \, t} =
\frac{1}{2 \, \pi} \, 2 \, W_{1} \, \frac{e^{j \,  W_{1} \, t}-e^{-j \, W_{1} \, t}}{2 \, j \, W_{1} \, t} =	
	\frac{W_{1}}{\pi} \frac{\sin\left( W_{1} \, t\right) }{W_{1} \, t}.
\end{equation}
Por lo tanto, resulta 
\begin{equation}
        x\left( t\right) =	\frac{W_{1}}{\pi} \frac{\sin\left( W_{1} \, t\right) }{W_{1} \, t}.
\end{equation}

\subsection{Exponencial decreciente derecha}

Consideremos la se\~{n}al 
\begin{equation}
	x\left( t\right) = e^{-a \, t}\, u\left( t\right) , a>0.
\end{equation}
Seg\'{u}n la ecuaci\'{o}n de an\'{a}lisis 
\begin{equation}
	X\left( j \, \omega \right) =
	\int_{0}^{\infty }e^{-a \, t} \, e^{-j \, \omega \, t} \, dt=
	\left[ -\frac{e^{-\left( a+j \, \omega \right) t}}{a+j \, \omega }\right] _{0}^{\infty },
\end{equation}
o sea que 
\begin{equation}
	X\left(j \, \omega \right) =\frac{1}{a+j \, \omega }, a>0.
\end{equation}
Observe adem\'{a}s que la siguiente integral est\'{a} acotada
\begin{equation}
	\int_{0}^{\infty }\left| e^{-a \, t}\, u\left( t\right) \right| dt=\frac{1}a.
\end{equation}
y adem\'{a}s la funci\'{o}n $x(t)$ tiene una \'{u}nica discontinuidad acotada en $t = 0$. 
Por lo tanto esta funci\'{o}n verifica las condiciones de Dirichlet.


\section{Propiedades de la Transformada de Fourier de tiempo continuo}

\subsection{Linealidad}

Sea una funci\'{o}n
\begin{equation}
        r\left( t\right) = a \, x\left( t\right) +b \, y\left( t\right).
\end{equation}
Podemos probar que su transformada es
\begin{equation}
        R\left( j \, \omega \right) =
        \int_{-\infty }^{\infty }\left( ax\left( t\right)+by\left( t\right) \right) e^{-j \, \omega \, t} \, dt.
\end{equation}
En otras palabras, la transformada de Fourier es una operaci\'{o}n lineal.

Escribiendo la ecuaci\'{o}n de an\'{a}lisis para $r\left( t\right)$
tenemos
\begin{equation}
        R\left( j \, \omega \right) =
        \left( \int_{-\infty }^{\infty } a \, x\left( t\right)e^{-j \, \omega \, t} + b \, y\left( t\right)e^{-j \, \omega \, t} \, dt\right),
\end{equation}
que puede rescribirse, por la linealidad del operador integral, como
\begin{equation}
        R\left( j \, \omega \right) =
        a \, \left( \int_{-\infty }^{\infty }x\left( t\right)e^{-j \, \omega \, t} \, dt\right) +
        b \, \left( \int_{-\infty }^{\infty }y\left( t\right)e^{-j \, \omega \, t} \, dt\right) .
\end{equation}
Por lo tanto resulta
\begin{equation}
        R\left( j \, \omega \right) =
                a \, X\left( j \, \omega \right) +
                b \, Y\left( j \, \omega \right).
\end{equation}
As\'{\i} deducimos la propiedad
\begin{equation}
        ax\left( t\right) +by\left( t\right) 
        \begin{array}{l}
                TF \\ 
                \leftrightarrow
        \end{array}
        a \, X\left( j \, \omega \right) +b \, Y\left( j \, \omega \right) .
\end{equation}

\subsection{Desplazamiento en el tiempo}

Dada la funci\'{o}n 
\begin{equation}
   r\left( t\right) = x\left( t-t_{0}\right) ,
\end{equation}
su transformada es
\begin{equation}
   R\left( j \, \omega \right) =  \int_{-\infty }^{\infty }x\left( t-t_{0}\right)e^{-j \, \omega \, t} \, dt.
\end{equation}

Escribiendo la ecuaci\'{o}n de an\'{a}lisis de $r\left( t\right)$ 
y aplicando un cambio de variable $\tau =t-t_{0},$ obtenemos
\begin{equation}
  R\left( j \, \omega \right) =
  \int_{-\infty }^{\infty }x\left( \tau \right) \, e^{-j \, \omega \left( \tau +t_{0}\right) }dt.
\end{equation}
Por las propiedades de la funci\'{o}n exponencial 
\begin{equation}
e^{-j \, \omega \, \left( \tau +t_{0}\right)} = e^{-j \, \omega \, \tau} \, e^{-j \, \omega \, t_{0}},
\end{equation}
reescribimos la ecuaci\'{o}n anterior como
\begin{equation}
   R\left( j \, \omega \right) = e^{-j \, \omega \, t_{0}} \int_{-\infty }^{\infty }x\left(\tau \right) e^{-j \, \omega \tau }d \tau.
\end{equation}
Observe que $e^{-j \, \omega \, t_{0}}$ sale de la integral pues no depende de la variable de integraci\'{o}n $\tau$.
As\'{'i}, llegamos a la conclusi\'{o}n 
\begin{equation}
        x\left( t-t_{0}\right) 
        \begin{array}{l}
                TF \\ 
                \leftrightarrow
        \end{array}
        e^{-j \, \omega \, t_{0}}X\left( j \, \omega \right) .
\end{equation}

\subsection{Conjugaci\'{o}n y simetr\'{\i}a del conjugado}

\begin{enumerate}
\item  El conjugado de la transformada de Fourier es 
\begin{equation}
      \overline{X\left( j \, \omega \right)} =
               \overline{\int_{-\infty }^{\infty }x\left(t\right) \, e^{-j \, \omega \, t} \, dt}=
               \int_{-\infty }^{\infty } \overline{x\left( t\right)} \, e^{j \, \omega \, t} \, dt.
\end{equation}
Cambiando el signo de $\omega$ en ambos lados de la ecuaci\'{o}n mantenemos la igualdad
\begin{equation}
        \overline{X\left( -j \, \omega \right)} =
                \int_{-\infty }^{\infty }\overline{x\left(t\right)} \, e^{-j \, \omega \, t} \, dt.
\end{equation}
Como conclusi\'{o}n, podemos escribir
\begin{equation}
        \overline{x\left( t\right)} 
        \begin{array}{l}
                TF \\ 
                \leftrightarrow
        \end{array}
        \overline{X\left( -j \, \omega \right)} .
\end{equation}

\item  Si $x\left( t\right) $ es real entonces 
\begin{equation}
        x\left( t\right) = \overline{x\left( t\right)}. \label{TFTCsimconj}
\end{equation}
Planteando la igualdad
\begin{equation}
        \overline{X\left( -j \, \omega \right)} =
                \overline{\int_{-\infty }^{\infty }x\left(t\right) \, e^{j \, \omega \, t} \, dt}=
                        \int_{-\infty }^{\infty }\overline{x\left( t\right)} \, e^{-j \, \omega \, t} \, dt,
\end{equation}
y considerando la Ec. (\ref{TFTCsimconj}) obtenemos
\begin{equation}
        \overline{X\left( -j \, \omega \right)} =
                \int_{-\infty }^{\infty }\overline{x\left(t\right)} e^{-j \, \omega \, t} \, dt=
                \int_{-\infty }^{\infty }x\left( t\right)e^{-j \, \omega \, t} \, dt=X\left( j \, \omega \right) ,
\end{equation}
lo que demuestra 
\begin{equation}
        X\left( -j \, \omega \right) = \overline{X\left( j \, \omega \right)} .  \label{conjuuno}
\end{equation}

\item  Si $x\left( t\right) $ es real por la ec. (\ref{conjuuno})
tenemos que 
\begin{equation}
\Re\left( X\left( j \, \omega \right) \right) =
\Re\left(X\left(-j \, \omega \right) \right),
\end{equation}
y 
\begin{equation}
 \Im\left( X\left( j \, \omega \right) \right) =
-\Im\left( X\left(-j \, \omega \right) \right).
\end{equation}

\item  Si $x\left( t\right) $ es real, entonces 
\begin{eqnarray}
\left\| X\left( j \, \omega \right) \right\| &=&
        X\left( j \, \omega \right) \, \overline{X\left( j \, \omega \right)} =
        X\left( j \, \omega \right) \, X\left( -j \, \omega \right) ,
\\
\left\| X\left( -j \, \omega \right) \right\| &=&
X\left( -j \, \omega \right) \, \overline{X\left( -j \, \omega \right)} =
X\left( -j \, \omega \right) \, X\left( j \, \omega \right) ,
\end{eqnarray}
debido a la ec. (\ref{conjuuno}), y por lo tanto $\left\| X\left( j \, \omega
\right) \right\| $ es una funci\'{o}n par, y de forma similar puede
demostrarse que $\angle X\left( j \, \omega \right) =\arctan \frac{%
\Im\left( X\left( j \, \omega \right) \right) }{\Re\left( X\left(
j \, \omega \right) \right) }$ es una funci\'{o}n impar.

\item  Si $x\left( t\right) $ es real y par entonces 
$X\left( j \, \omega\right) $ es real y par pues 
\begin{equation}
X\left( -j \, \omega \right) =\int_{-\infty }^{\infty }x\left( t\right)
e^{j \, \omega \, t} \, dt=
\int_{-\infty }^{\infty }x\left( -t \right) e^{-j \, \omega \, t }dt .
\end{equation}
Como $x\left( t\right) = x\left( -t\right) $ (funci\'{o}n par) tenemos
que 
\begin{equation}
X\left( -j \, \omega \right) =\int_{-\infty }^{\infty }x\left( \tau \right)
e^{-j \, \omega \tau }d\tau =X\left( j \, \omega \right) ,
\end{equation}
y por lo tanto es una funci\'{o}n par.

\item  Si $x\left( t\right) $ es real e impar, $\left( x\left( t\right)
=-x\left( -t\right) \right) ,$ entonces $X\left( j \, \omega \right) $ es impar
pues 
\begin{equation}
-X\left( -j \, \omega \right) =-\int_{-\infty }^{\infty }x\left( t\right)
e^{j \, \omega \, t} \, dt=\int_{-\infty }^{\infty }-x\left( -\tau \right) e^{-j \, \omega
\tau }d\tau ,
\end{equation}
\begin{equation}
-X\left( -j \, \omega \right) =\int_{-\infty }^{\infty }-x\left( -\tau \right)
e^{-j \, \omega \tau }d\tau =\int_{-\infty }^{\infty }x\left( \tau \right)
e^{-j \, \omega \tau }d\tau =X\left( j \, \omega \right) .
\end{equation}
y considerando la ec. (\ref{conjuuno}), adem\'{a}s $X\left( j \, \omega \right) $
es imaginaria pura 
\begin{equation}
X\left( j \, \omega \right) =-X\left( -j \, \omega \right).
\end{equation}
\end{enumerate}

\subsection{Diferenciaci\'{o}n en tiempo}

\begin{enumerate}
\item Escribamos la ecuaci\'{o}n de s\'{\i}ntesis
\begin{equation}
x\left( t\right) =\frac{1}{2 \, \pi }\int_{-\infty }^{\infty }X\left( j \, \omega
\right) e^{j \, \omega \, t}d\omega .
\end{equation}

\item Derivemos ambos t\'{e}rminos respecto al tiempo
\begin{equation}
\frac{dx\left( t\right) }{dt}=\frac{1}{2 \, \pi }\int_{-\infty }^{\infty
}X\left( j \, \omega \right) \frac{d\left( e^{j \, \omega \, t}\right) }{dt}d\omega.
\end{equation}

\item Obtenemos finalmente
\begin{equation}
\frac{dx\left( t\right) }{dt}= \frac{1}{2 \, \pi }\int_{-\infty
}^{\infty }j \, \omega \, X\left( j \, \omega \right) e^{j \, \omega \, t}d\omega.
\end{equation}

\item La conclusi\'{o}n es 
\begin{equation}
\frac{dx\left( t\right) }{dt} 
\begin{array}{l}
F \\ 
\leftrightarrow
\end{array}
j \, \omega \, X\left( j \, \omega \right) .
\end{equation}

\end{enumerate}

\subsection{Diferenciaci\'{o}n en frecuencia}

\begin{enumerate}
\item Escribamos la ecuaci\'{o}n de an\'{a}lisis
\begin{equation}
X\left( j \, \omega \right) =\int_{-\infty }^{\infty }x\left( t\right)
e^{-j \, \omega \, t} \, dt.
\end{equation}

\item Derivemos ambos t\'{e}rminos respecto a la frecuencia
\begin{equation}
\frac{dX\left( j \, \omega \right) }{d\omega }=\int_{-\infty }^{\infty }x\left(
t\right) \, \frac{d\left( e^{-j \, \omega \, t}\right) }{d\omega } \, dt.
\end{equation}

\item Obtenemos
\begin{equation}
\frac{dX\left( j \, \omega \right) }{d\omega }=
\int_{-\infty }^{\infty}-j \, t \, x\left( t\right) \, e^{-j \, \omega \, t} {d\omega } \, dt.
\end{equation}

\item Finalmente 
\begin{equation}
j \, \frac{dX\left( j \, \omega \right) }{d\omega }=\int_{-\infty }^{\infty
}t \, x\left( t\right) \, e^{-j \, \omega \, t} {d\omega } \, dt.
\end{equation}

\end{enumerate}

Conclusi\'{o}n 
\begin{equation}
t \, x\left( t\right) 
\begin{array}{l}
F \\ 
\leftrightarrow
\end{array}
j \, \frac{dX\left( j \, \omega \right) }{d\omega }.
\end{equation}




\subsection{Escalado en tiempo y frecuencia}

\begin{enumerate}
	
\item Vamos a calcular la transformada de $x(a \, t)$ en función de la transformada de $x(t)$.
 
\item Escribamos la ecuaci\'{o}n de an\'{a}lisis
\begin{equation}
F\left( x\left( at\right) \right) =\int_{-\infty }^{\infty }x\left(
a \, t\right) e^{-j \, \omega \, t} \, dt,
\end{equation}
donde $a\neq 0.$ 

\item Hagamos la sustituci\'{o}n $\tau =a \, t$ para obtener
\begin{equation}
F\left( x\left( at\right) \right) =\left\{ 
\begin{array}{c}
\frac{1}a\int_{-\infty }^{\infty }x\left( \tau \right) e^{-\frac{j \, \omega 
}a\tau }d\tau , a>0, \\ 
\\ 
-\frac{1}a\int_{-\infty }^{\infty }x\left( \tau \right) e^{-\frac{j \, \omega 
}a\tau }d\tau , a<0.
\end{array}
\right.
\end{equation}

\item El signo de $a$ afecta al orden de integraci\'{o}n. Si se conservan los
l\'{\i}mites de integraci\'{o}n, debe observarse el signo correspondiente. La
conclusi\'{o}n es 
\begin{equation}
x\left( at\right) 
\begin{array}{l}
F \\ 
\leftrightarrow
\end{array}
\frac{1}{\left| a\right| }X\left( \frac{j \, \omega }a\right) ,
\end{equation}
donde $a\neq 0.$

\end{enumerate}

\subsection{Convoluci\'{o}n en el dominio del tiempo}

\begin{enumerate}
\item Empezamos con la integral de convoluci\'{o}n 
\begin{equation}
        y\left( t\right) =
                \int_{-\infty }^{\infty }       x\left( \tau \right) \, h\left(t-\tau \right)      \, d\tau .
\end{equation}
\item Deseamos obtener la transformada de Fourier 
\begin{equation}
        Y\left( j \, \omega \right) =
                F\left( y\left( t\right) \right) =
                \int_{-\infty}^{\infty }
                  \left[ \int_{-\infty }^{\infty }x\left( \tau \right) \, h\left(t-\tau \right) d\tau \right] \,
                   e^{-j \, \omega \, t} \, dt.
\end{equation}
\item Cambiando el orden de integraci\'{o}n, y observando que $x\left( \tau\right) $ no depende de $t,$ obtenemos 
\begin{equation}
        \int_{-\infty }^{\infty }
                \left[ \int_{-\infty }^{\infty }
                        x\left( \tau\right) \, h\left( t-\tau \right) 
                d\tau \right] \, e^{-j \, \omega \, t}
        dt =
        \int_{-\infty}^{\infty }
                x\left( \tau \right) 
                \left( \int_{-\infty }^{\infty } \, h\left(t-\tau \right) \, e^{-j \, \omega \, t}    \, dt\right) 
        d\tau .
\end{equation}
\item Por la propiedad de desplazamiento en el tiempo, aplicada a $h(t - \tau)$, tenemos que 
\begin{equation}
        \int_{-\infty }^{\infty }h\left( t-\tau \right) \, e^{-j \, \omega \, t} \, dt=
                e^{-j \, \omega \tau } \, h\left( j \, \omega \right) ,
\end{equation}
que podemos reemplazar para obtener 
\begin{equation}
        Y\left( j \, \omega \right) =
        \int_{-\infty }^{\infty }
                x\left( \tau \right)
                \left( \int_{-\infty }^{\infty }h\left( t-\tau \right) \, e^{-j \, \omega \, t} \, dt\right) 
                d\tau =
       \int_{-\infty }^{\infty }x\left( \tau \right) \, e^{-j \, \omega\tau } \, H\left( j \, \omega \right) d\tau .
\end{equation}
\item La integral entre par\'{e}ntesis 
\begin{equation}
        Y\left( j \, \omega \right) =
        \int_{-\infty }^{\infty }
        \left( x\left( \tau\right) e^{-j \, \omega \tau } \, d\tau \right) \, H\left( j \, \omega \right),
\end{equation}
es igual a la transformada de Fourier $X\left( j \, \omega \right) $ y por lo
tanto 
\begin{equation}
        Y\left( j \, \omega \right) = X\left( j \, \omega \right) \, H\left( j \, \omega \right) .
\end{equation}

\item Conclusi\'{o}n 
\begin{equation}
y\left( t\right) =x\left( t\right) * h\left( t\right) 
        \begin{array}{l}
                F \\ 
                \leftrightarrow
        \end{array}
        Y\left( j \, \omega \right) = X\left( j \, \omega \right) R\left( j \, \omega \right) .
\end{equation}

\end{enumerate}

\textbf{Convoluci\'{o}n de se\~{n}ales en el dominio del tiempo implica multiplicaci\'{o}n de 
transformadas de Fourier en el dominio de la frecuencia}

\subsection{Convolución en el dominio de la frecuencia}

\begin{enumerate}
\item Empezamos con la expresi\'{o}n 
\begin{equation}
        R\left( j \, \omega \right) =
        \frac{1}{2 \, \pi }
        \int_{-\infty }^{\infty }
             X\left(j \, \omega \right) Y\left( j\left( \omega -\theta \right) \right) d\theta .
\end{equation}

\item Aplicamos la ecuaci\'{o}n de s\'{\i}ntesis 
\begin{equation}
        r\left( t\right) =
                \frac{1}{2 \, \pi }
                \int_{-\infty }^{\infty }
                \left[ \frac{1}{2 \, \pi }
                          \int_{-\infty }^{\infty }
                              X\left( j \, \omega \right) \, Y\left( j\left(\omega -\theta \right) \right) \, 
                          d\theta 
                \right] \, e^{j \, \omega \, t} \,
                d\omega .
\end{equation}

\item Intercambiamos el orden de integraci\'{o}n 
\begin{equation}
        r\left( t\right) =
        \frac{1}{2 \, \pi }
        \left[ \int_{-\infty }^{\infty }
                X\left(j \, \omega \right) \left( \frac{1}{2 \, \pi }
                \int_{-\infty }^{\infty }
                    Y\left(j\left( \omega -\theta \right) \right) e^{j \, \omega \, t}d\omega \right) d\theta 
        \right].
\end{equation}

\item Aplicamos la ecuaci\'{o}n de s\'{\i}ntesis 
\begin{equation}
        r\left( t\right) =
                \frac{1}{2 \, \pi }
                \left[ 
                     \int_{-\infty }^{\infty }X\left(j \, \omega \right) \, y\left( t\right) \, e^{j\theta t} \, d\theta 
                \right].
\end{equation}
y dado que $y\left( t \right) $ no depende de la variable de integraci\'{o}n lo podemos extraer de la 
integral 
\begin{equation}
   r\left( t\right) =
       y\left( t\right) 
      \frac{1}{2 \, \pi } \left[ \int_{-\infty}^{\infty }X\left( j \, \omega \right) \,  e^{j \, \theta \, t} \, d\theta \right] .
\end{equation}

\item Finalmente, aplicando nuevamente la ecuaci\'{o}n de s\'{\i}ntesis, obtenemos
\begin{equation}
        r\left( t\right) = y\left( t\right) x\left( t\right) .
\end{equation}

\end{enumerate}

Conclusi\'{o}n
\begin{equation}
        r\left( t\right) =x\left( t\right) \, y\left( t\right) 
        \begin{array}{l}
                F \\ 
                \leftrightarrow
        \end{array}     
        R\left( j \, \omega \right) =
                \frac{1}{2 \, \pi }
                        \int_{-\infty }^{\infty }
                                X\left(j \, \omega \right) \, Y\left( j\left( \omega -\theta \right) \right) \,
                        d\theta.
\end{equation}
\textbf{Convoluci\'{o}n de transformadas implica multiplicaci\'{o}n de se\~{n}ales en el dominio del 
tiempo}

\subsection{Traslaci\'{o}n en frecuencia}

Empezamos con la integral de s\'{\i}ntesis

\begin{equation}
\begin{array}{l}
     \frac{1}{2 \, \pi }\int_{-\infty }^{\infty }
         X\left( j\left( \omega -\omega_{0}\right) \right) e^{j \, \omega \, t}d\omega = \\ 
        \\ 
        \frac{e^{j \, \omega _{0}t}}{2 \, \pi } \int_{-\infty }^{\infty }
                X\left( j\left(\omega -\omega _{0}\right) \right) e^{j\left( \omega -\omega _{0}\right)t}
                d\left( \omega -\omega _{0}\right) ,
\end{array}
\end{equation}

\begin{equation}
        \frac{e^{j \, \omega_{0}t}}{2 \, \pi }
        \int_{-\infty }^{\infty }X\left( j \, \omega_{1}\right) e^{j \, \omega _{1}t}d\omega _{1}=
        e^{j \, \omega _{0}t}x\left( t\right) .
\end{equation}

Conclusi\'{o}n 
\begin{equation}
e^{j \, \omega _{0}t}x\left( t\right)  
\begin{array}{l}
F \\ 
\leftrightarrow
\end{array}
 X\left( j\left( \omega -\omega _{0}\right) \right) .
\end{equation}
Recordemos la propiedad de desplazamiento temporal para comparar 
\begin{equation}
x\left( t-t_{0}\right)  
\begin{array}{l}
F \\ 
\leftrightarrow
\end{array}
 e^{-j \, \omega \, t_{0}}X\left( j \, \omega \right) .
\end{equation}

\subsection{Relaci\'{o}n de Parseval}

Empezamos con la igualdad
\begin{equation}
\int_{-\infty }^{\infty }
  \left| x\left( t\right) \right|^{2} \, dt=
  \int_{-\infty }^{\infty }x\left( t\right) \, \overline{x^{ }\left(
      t\right)} \, dt.
\end{equation}
Recordamos la ecuaci\'{o}n de s\'{\i}ntesis 
\begin{equation}
\int_{-\infty }^{\infty }x\left( t\right) \, \overline{x\left( t\right)}
\, dt=
  \int_{-\infty }^{\infty }
    x\left( t\right) 
    \left[\frac{1}{2 \, \pi }\int_{-\infty }^{\infty }\overline{X\left(
          j \, \omega \right)} e^{-j \, \omega \, t} \, d\omega \right]
  dt,
\end{equation}
y cambiamos el orden de integraci\'{o}n 
\begin{equation}
\int_{-\infty }^{\infty }x\left( t\right) \, \overline{x\left(
    t\right)} \, dt=
\frac{1}{2 \, \pi }\int_{-\infty }^{\infty }\overline{X\left( j \, \omega
  \right)} \, \left(
\int_{-\infty }^{\infty }x\left( t\right) e^{-j \, \omega \, t} \, dt\right) d\omega .
\end{equation}
Recordemos la ecuaci\'{o}n de an\'{a}lisis y escribimos 
\begin{equation}
\int_{-\infty }^{\infty }x\left( t\right) \, \overline{x\left(
    t\right)} \, dt=
\frac{1}{2 \, \pi }\int_{-\infty }^{\infty }\overline{X\left( j \, \omega
  \right)} \, X\left(j \, \omega \right) \, d\omega .
\end{equation}
Conclusi\'{o}n: si es limitada, la energ\'{\i}a de una se\~{n}al puede calcularse en el dominio del tiempo
y en el dominio de la frecuencia
\begin{equation}
        \int_{-\infty }^{\infty }\left| x\left( t\right) \right| ^{2}dt=
        \frac{1}{2 \, \pi }\int_{-\infty }^{\infty }\left| X\left( j \, \omega \right) \right|^{2}d\omega .
\end{equation}

\subsection{Dualidad}

La demostraci\'{o}n de la propiedad de dualidad es m\'{a}s sencilla si consideramos las ecuaciones de 
an\'{a}lisis y s\'{\i}ntesis en la variable $f$
\begin{equation}
  X_{1}\left( f\right) =\int_{-\infty }^{\infty }x_{1}\left( t\right) e^{-j2 \, \pi ft}dt,    \label{anaFouDua}
\end{equation}
y 
\begin{equation}
  x_{1}\left( t\right) =\int_{-\infty }^{\infty }X_{1}\left( f\right) e^{j2 \, \pi ft}df.     \label{sinFouDua}
\end{equation}

Observemos que las dos ecuaciones son muy similares. Reemplazemos ahora $t = -\tau$ para obtener
\begin{equation}
      X_{2}\left( f\right) = \int_{-\infty }^{\infty }x_{2}\left( -\tau\right) e^{j2 \, \pi f\tau}d(-\tau),
       \label{duaAnaFou}
\end{equation}
y 
\begin{equation}
      x_{2}\left(-\tau\right) =\int_{-\infty }^{\infty }X_{2}\left( f\right) e^{-j2 \, \pi f \tau}df.
       \label{duaSinFou}
\end{equation}

Comparemos los pares de ecuaciones (\ref{anaFouDua}) y (\ref{duaSinFou}), (\ref{sinFouDua}) y 
(\ref{duaAnaFou}). Podemos ver que tienen la misma forma, m\'{a}s alla de que la variable de 
integraci\'{o}n es distinta. De esta similitud entre las ecuaciones es posible deducir la relaci\'{o}n de 
dualidad
\begin{equation}
        \begin{array}{cc}
                x(t)    \begin{array}{l}
                                F \\ 
                                \leftrightarrow
                        \end{array}
                X(f)    \\
                X(-t)   \begin{array}{l}
                                F \\ 
                                \leftrightarrow
                        \end{array}
                x(f)
        \end{array}
\end{equation}

\subsection{Integraci\'{o}n}

Dada una se\~{n}al $x(t)$, la integral de $x(t)$ es la funci\'{o}n $\int_{-\infty}^{t} x(v) dv$.
Supongamos que $x(t)$ tiene una transformada de Fourier $X(\omega)$. No podemos asegurar que la
integral de $x(t)$ tenga una transformada de Fourier en el sentido ordinario, pero tiene una
transformada generalizada. Se puede demostrar que 
\begin{equation}
\int_{-\infty }^{t}x\left( t\right) dt 
\begin{array}{l}
F \\ 
\leftrightarrow
\end{array}
\frac{1}{j \, \omega }X\left( j \, \omega \right) +\pi X\left( 0\right) \delta\left( \omega \right),
\end{equation}
donde $\delta(\omega)$ es un impulso en el dominio de la frecuencia. El t\'{e}rmino 
$\pi X\left( 0\right) \delta\left( \omega \right)$ aparece para tomar en cuenta el valor 
promedio de la integral.

Para probar esta propiedad, debemos encontrar primero la transformada del impulso $\delta(t)$
y del escal\'{o}n $u(t)$.

\subsubsection{Impulso $\delta(t)$}

\begin{enumerate}

\item Por la propiedad del muestreo
\begin{equation}
  \delta(t) e^{j \omega \, t} = \delta(t).
\end{equation}

\item Por lo tanto, 
\begin{equation}
  \delta(t) 
  \begin{array}{l}
    F \\ 
  \leftrightarrow
  \end{array}
  \int_{-\infty}^{\infty} \delta(t) dt = 1.
\end{equation}

\end{enumerate}

\subsubsection{Escal\'{o}n $u(t)$}

\begin{enumerate}

\item Primero encontremos la transformada de 
\begin{equation}
x(t) = -\frac{1}{2} + u(t).
\end{equation}

\item Dado que la se\~{n}al $x(t)$ es constante excepto por un salto de valor $1$ 
desde $t=0^{-}$ a $t=0^{+}$, la derivada (generalizada) de $x(t)$ es igual al 
impulso unitario $\delta(t)$. Entonces por la propiedad de la derivada, la transformada
de Fourier de $x(t)$ debe ser
\begin{equation}
  x(t) = \frac{1}{2} + u(t)
  \begin{array}{l}
    F \\ 
  \leftrightarrow
  \end{array}
  \frac{1}{j \omega}. \label{fourierintegraluno}
\end{equation}

\item Recordemos que la transformada de Fourier de la se\~{n}al constante $\frac{1}{2}$ es
$\pi \delta(\omega)$. Empleando la ec. (\ref{fourierintegraluno}) y la propiedad de linealidad,
obtenemos finalmente
\begin{equation}
  u(t) 
  \begin{array}{l}
    F \\ 
  \leftrightarrow
  \end{array}
  \pi \delta(\omega) + \frac{1}{j \omega}.
\end{equation}

\end{enumerate}

\subsubsection{Integral}

\begin{enumerate}

\item La integral de $x(t)$ puede ser expresada como
\begin{equation}
  \int_{-\infty}^{t} x(v) dv = u(t) * x(t).
\end{equation}
Entonces por la propiedad de la convoluci\'{o}n, la transformada de Fourier de la integral 
$x(t)$ es igual a $U(\omega) X(\omega)$.

\item Por lo tanto
\begin{equation}
\int_{-\infty }^{t}x\left( t\right) dt 
\begin{array}{l}
F \\ 
\leftrightarrow
\end{array}
(\frac{1}{j \, \omega } +\pi \delta(\omega) ) X\left( \omega \right) = 
\frac{X(\omega)}{j \omega} + \pi X(0) \delta(\omega).
\end{equation}

\end{enumerate}
